\section{Current Situation}

The destruction of cultural heritage in armed conflict is not a relic of history. It is happening right now and is accelerating. The early 21st Century has produced a sustained, global assault on the places that human communities use to remember who they are. What makes the current moment distinct is not simply the scale of destruction, it is the growing body of evidence that much of this destruction is not accidental. It is deliberate, strategic, and, in many cases, goes entirely unpunished.

\subsection{Heritage Destruction as a Strategy of War}

Despite decades of legal frameworks insisting otherwise, the world still does not fully treat the destruction of cultural heritage as the weapon it frequently is. Heritage destruction functions as a weapon of war in several distinct and overlapping ways.

The most visible is identity erasure. Every community anchors its sense of self in physical place. Destroying these places does not merely damage property. It severs the connection between a living community and its past, sending a message that is simultaneously military and existential. The second function is demoralization. Populations under attack draw psychological sustenance from the continuity of familiar places. Military strategists have long understood this and utilised this as part of their strategies. Third, heritage destruction is frequently used to eliminate evidence. Archaeological sites contain physical proof of habitation  proof that a particular people lived in a particular place for a particular length of time. Destroying those sites destroys the evidence. This is not incidental: in conflicts where competing territorial claims rest partly on historical presence, the elimination of archaeological proof is a strategic act with long-term political consequences.

Finally, heritage destruction serves as an instrument of ideological assertion. Religious symbols are often destroyed with the intent of not only erasing its existence in that area, but to leave space for a different religion, ethnicity, or nation to take its place.

What all of these functions share is intentionality. They are not accidents of war. They are choices  and the failure to treat them as such has allowed perpetrators to hide behind the fog of conflict, claiming that damage was unavoidable, that military necessity prevailed, that the real target was something else. The documentation work being done by UNESCO, by satellite imagery analysts, by investigative journalists and by local community members is building an evidentiary record that increasingly makes these excuses impossible to sustain.

\subsection{State Responsibility and Non-State Armed Groups}

The central accountability gap in international heritage law can be stated simply: the framework was designed for a world of states fighting each other, and that is not the world we live in.

The 1954 Hague Convention binds states. The Rome Statute can prosecute individuals. What neither instrument fully addresses is the sprawling middle ground; the non-state armed groups that have carried out some of the most systematic heritage destruction of the modern era, and the states that arm, fund, or tolerate them without ever formally directing their actions.

When state armed forces destroy cultural heritage, the legal picture is theoretically clearer. A state whose military deliberately targets a protected site is in violation of the Hague Convention. Its individual commanders may be liable under the Rome Statute for war crimes. The state itself may bear responsibility under customary international law for internationally wrongful acts. In practice, however, accountability for state actors remains almost entirely theoretical. No state has ever been formally convicted of heritage destruction as a distinct internationally wrongful act.

The challenge posed by non-state armed groups is different and in some ways even more intractable. Groups like ISIS, Ansar Dine, the Taliban, and various militia forces operating across the Sahel, the Middle East, and beyond are not parties to the Hague Convention. They cannot be sued before international courts in the traditional sense. Their individual members may theoretically be prosecuted under the Rome Statute; the Al Mahdi case proved that this is possible but prosecution requires arrest, arrest requires custody, and custody requires either the group's defeat or extraordinary international cooperation which rarely materializes.

This leaves the question of state complicity. When a state funds, arms, or provides political cover to a non-state group that goes on to destroy cultural heritage, what responsibility does that state bear? International law has wrestled with this question in the context of human rights violations for decades, and the answers remain contested. The degree of control a state must exercise over a group before its actions become attributable to that state the so-called "effective control" test established in the Nicaragua case  sets a very high bar. In practice, states can provide substantial support to non-state actors engaged in heritage destruction while maintaining plausible deniability.

There is also the question of omission. States are obliged under international law not merely to refrain from destroying cultural heritage themselves but to take measures to prevent its destruction by others within their territory or control. When a government is unable or unwilling to protect heritage sites under threat from armed groups operating within its borders, the question of responsibility becomes deeply uncomfortable particularly when the international community is simultaneously denying that government the resources or military support it would need to actually protect those sites.

The current legal framework creates a perverse situation: the actors most responsible for heritage destruction are often those least affected by accountability mechanisms, while the actors most reachable states can deflect responsibility through the non-state groups they enable.

\subsection{Challenges to Accountability}

Even where the legal framework theoretically applies and political will theoretically exists, the practical obstacles to accountability for heritage destruction are formidable.

Documentation during active conflict is the most immediate challenge. Evidence of heritage destruction must be gathered as quickly as possible after damage occurs, before further destruction, looting, or reconstruction obscures the record. Active conflict zones are precisely the places where access is most restricted, where local experts are most at risk, and where the infrastructure for systematic documentation is the most likely to have collapsed. UNESCO has developed remote-sensing monitoring tools and partnerships with satellite imagery providers to track damage from a distance, but while remote sensing can identify that destruction has occurred; it cannot always establish how, by whom, or with what intent. The difference between targeted destruction and collateral damage is a legal and moral distinction of enormous consequence, which is also the hardest to prove.

Establishing intent is the second major challenge. International law does not criminalize all wartime damage to cultural heritage, only deliberate attacks. Proving deliberateness requires demonstrating that the attacker knew the site had cultural significance, that the site was not being used for military purposes at the time of the attack, and that the attacker chose to strike it anyway. Each of these elements is difficult to establish, particularly when the perpetrating party controls the relevant communications, intelligence assessments, and military records. Perpetrators have become increasingly sophisticated in constructing alternative narratives: cultural sites were being used as military positions; the destruction was an unavoidable consequence of legitimate military action; the targeting was carried out by the other side. Countering these narratives requires the kind of detailed, meticulous, contemporaneous documentation that is extremely hard to produce under fire.

Jurisdictional limitations compound these difficulties. The International Criminal Court can only prosecute individuals who are nationals of states which have ratified the Rome Statute, or whose crimes were committed on the territory of such states  unless the UN Security Council refers the situation to the Court. This means that nationals of non-signatory states, including major military powers, are effectively beyond the Court's reach unless the Security Council acts  which, as noted, the veto structure makes it extremely unlikely when a permanent member is involved. The Al Mahdi conviction was possible because Mali had ratified the Rome Statute, and because the perpetrators were members of a non-state group without a powerful patron to block proceedings.

The tension between accountability and reconstruction is a challenge that is rarely acknowledged but deserves serious attention. The international community's instinct after heritage destruction is often to rebuild  to restore the lost monument, reconstruct the damaged mosque, and replicate the destroyed library. This impulse is understandable and, in many cases, genuinely valuable for community healing. The reconstruction of Timbuktu's mausoleums, undertaken with UNESCO support after the Al Mahdi conviction, was meaningful precisely because it accompanied rather than replaced a process of legal accountability.

Political selectivity may be the deepest challenge of all. The international response to heritage destruction has never been consistent. The destruction of European heritage often generates immediate, sustained, high-profile international attention. The destruction of heritage in non-Western states generates condemnations and then, often, silence. This double standard is not merely an ethical failure; it is a practical obstacle to building the kind of universal, principled accountability framework that would actually deter future destruction. As long as accountability is selectively applied based on geopolitical interest rather than consistent legal principle, perpetrators everywhere will correctly calculate that their chances of facing consequences depend less on what they did than on who they are and who their friends are.

This is the landscape delegates must navigate. The challenge is not simply to condemn heritage destruction; the international community has been doing that for decades without much effect. The challenge is to design accountability mechanisms that can actually reach the perpetrators, actually function during active conflict rather than only in its aftermath, and actually apply with enough consistency to function as a genuine deterrent. That is a harder, more specific, and more urgent task.
